\scanpage{30}

qui sont $\varprojlim$ du type envisagé de foncteurs représentables (dans le cas resp., on utilise, pour \deleted{\unsure{y}} donner un sens à cet énoncé -- dans $\widehat A_{\mathrm{ab}}$ -- la structure \textquotedblleft$\frac12$ additive\textquotedblright\ sur $A$ héritée par son plongement dans $C$ additive). Lorsque de plus $C$ est stable pour ce type de $\varprojlim$ (p. ex. des produits finis), alors on a \unsure{$=$|même} une \uline{équivalence} entre $C$ et la cat. des foncteurs $A^\circ\mathrel{-}\mathrm{Ens}$ (resp $A^\circ\mathrel{-}\mathrm{Ab}$) qui sont de telles $\varprojlim$ de foncteurs représentables.

\deleted{[Si par ex. $A$ est réduite à un objet, on trouve dans le cas $C$ additive] \unsure{cas resp., si tt objet|resp., si $A$ objet}}

Si p. ex. $C$ est additive, donc \deleted{$A$} induit sur $A$ une structure $\frac12$ additive, et si tt objet de $C$ est isomorphe à une somme directe \inserted{finie} d'objets de $A$, alors $C\simeq$ sous-cat. pleine de $\widehat A_{\mathrm{ab}}$ (foncteurs additifs contravariants de $A$ dans $\mathrm{Ab}$) formée \deleted{sont} des foncteurs isom. à une somme directe finie de foncteurs repr.

% archive sheet number: 29
